\section{Розділ~\ref{sec:eetypes}. Нейрони як прилади}

Режими роботи окремих клітин виміряно прямо, струмом через електрод і записом напруги; математика інтегратора й поділу на класи --- класична теорія, а те, що мозок користується переналаштуванням режимів для обчислень, лишається гіпотезою.

{\footnotesize
\setlength{\tabcolsep}{3pt}\renewcommand{\arraystretch}{1.15}
\begin{longtable}{>{\raggedright}p{0.5cm}>{\raggedright}p{4.0cm}>{\raggedright}p{5.6cm}>{\raggedright}p{2.3cm}>{\raggedright\arraybackslash}p{1.7cm}}
\toprule
№ & Твердження & Дослід і що виміряно & Джерело & Статус \\
\midrule\endfirsthead
\toprule
№ & Твердження & Дослід і що виміряно & Джерело & Статус \\
\midrule\endhead
1 & Резонатор: повільний струм, що протидіє зміні напруги, робить нейрон смуговим фільтром із максимумом на одиницях--десятках герц (підрозділ~\ref{sec:resonator}) & У зрізах у \nt{GLUT}-нейрони вводили синусоїдальний струм із плавно зростаючою частотою і вимірювали імпеданс $|Z(f)|$: максимум на $2$--$5$~Гц за $33^\circ$C (близько $7$~Гц за $38^\circ$C); резонанс створюють повільні калієвий і катіонний струми, а підсилює постійний натрієвий струм. Огляд дає той самий прийом для багатьох класів клітин & \cite{HuVervaekeStorm2002,HutcheonYarom2000} & виміряно \\
2 & Повільний струм поводиться як індуктивність, разом з ємністю мембрани --- коливальний контур & Підпорогову відповідь аксона на струм виміряно й описано лінеаризованими рівняннями провідностей; еквівалентна схема містить «феноменологічну» індуктивність, величина якої залежить від напруги & \cite{Mauro1970} & теорія \\
3 & Стала часу групи зі зворотним зв'язком $\tau_{\text{еф}}=\tau/(1-w)$~\eqref{eq:perfectint}; за $\tau=0.1$~с і $w=0.995$ --- $20$~с & Виводиться в розділі з лінійного рівняння групи; як механізм утримання положення ока запропоновано в теоретичній роботі & виведення в розділі; \cite{Seung1996} & теорія \\
4 & Інтегратор положення ока тримає вхід мережею, а не властивістю окремої клітини & Внутрішньоклітинний запис у риби у вузлі, що тримає положення ока: після швидкого повороту частота нейрона сходинкою змінюється й тримається; короткий струм в одну клітину спричиняє лише короткочасний зсув, а сходинки напруги зберігаються й нижче порога --- їх підтримують синаптичні входи & \cite{Aksay2001} & виміряно \\
5 & Інтеграторові без витоку потрібне постійне підлаштування навчанням; коли його розладнано, він або забуває, або йде в насичення & Рибу навчали зоровим зворотним зв'язком, пропорційним $\pm$ положенню ока: утримання ставало нестійким або з витоком, стала часу нейронів падала більш ніж на порядок, до $<1$~с; звичайний зоровий зворотний зв'язок поступово повертав стійкість & \cite{Major2004} & виміряно \\
6 & Бістабільний нейрон: короткий вхід умикає тривале плато, гальмування вимикає; властивість умикається серотоніном (підрозділ~\ref{sec:bistable}) & Внутрішньоклітинний запис \nt{ACET}-нейронів, що керують м'язами, у кішки: коротке синаптичне збудження переводить нейрон у тривалу роботу, коротке гальмування скидає; у спінальної кішки стан з'являється після введення попередника серотоніну & \cite{Hounsgaard1988} & виміряно \\
7 & Два класи: інтегратори (частота зростає від нуля) і резонатори (на порозі одразу скінченна частота) & Класи виведено з теорії біфуркацій. Дослід: у зрізах кори \nt{GLUT}-нейрони починають працювати з як завгодно малою частотою (тип~1), швидкі \nt{GABA}-нейрони --- стрибком зі скінченної частоти (тип~2) & \cite{Izhikevich2007,Tateno2004} & виміряно \\
8 & Один нейрон --- інтегратор або детектор збігів: гальмування вкорочує вікно додавання & У зрізах гіпокампа вікно, у якому збуджувальні входи додаються до порога, коротше за $2$~мс; його вкорочує пряме гальмування, що надходить одразу за збудженням; у дендритах, де гальмування слабше, вікно ширше & \cite{Pouille2001} & виміряно \\
9 & Лінія затримки з аксонів (таблиця~\ref{tab:eetypes}) & У сипухи виміряно затримки в аксонах, що йдуть до детекторів збігів: різна довжина аксона дає різну затримку, і детектори налаштовані на різницю часів надходження звуку & \cite{Carr1990} & виміряно \\
10 & Ритмоводій під час неспання й мовчазна клітина уві сні & \nt{SERO}-нейрони працюють майже регулярно вдень, сповільнюються в повільному сні й майже мовчать у швидкому (докладніше --- розділ~\ref{sec:sleep}) & \cite{JacobsAzmitia1992,McGintyHarper1976} & виміряно \\
11 & Прилад задають іонні канали й широкомовні канали, що їх підлаштовують (твердження~\ref{prop:eetypes}) & Показано на малих мережах: речовини-регулятори змінюють власні властивості нейронів і силу синапсів, тож та сама схема видає різні ритми & \cite{Marder2012} & виміряно \\
12 & Мозок використовує переналаштування режимів для обчислень (твердження~\ref{prop:eetypes}) & Прямого досліду, де переналаштування режиму клітини було б потрібне для розв'язання задачі, немає; є лише досліди, де регулятори змінюють вихід схеми & \cite{Marder2012} & гіпотеза \\
\bottomrule
\end{longtable}}
